\subsection{INTEGRATION OF COMPARISON RELATIONS: II. STRONG RELATIONS}

PROPOSITION 6. Let \(\mathbf{f}\) be a regulated vector function, and \(g\) a regulated function \(\geqslant 0\) on \([a, +\infty[\) .

\(1^{\circ}\) If the integral \(\int_{a}^{+\infty} g(t) dt\) converges then the relation \(\mathbf{f} \ll g\) (resp. \(\mathbf{f} \sim \mathbf{c}g\) , where \(\mathbf{c}\) is constant) implies that \(\int_{x}^{+\infty} \mathbf{f}(t) dt \ll \int_{x}^{+\infty} g(t) dt\) (resp. \(\int_{x}^{+\infty} \mathbf{f}(t) dt \sim \mathbf{c} \int_{x}^{+\infty} g(t) dt\) ).

\(2^{\circ}\) If the integral \(\int_{a}^{+\infty} g(t) dt\) is infinite then the relation \(\mathbf{f} \ll g\) (resp. \(\mathbf{f} \sim \mathbf{c}g\) ) implies that

\[
\int_ {\alpha} ^ {x} \mathbf {f} (t) d t \ll \int_ {\beta} ^ {x} g (t) d t \quad (\text { resp. } \quad \int_ {\alpha} ^ {x} \mathbf {f} (t) d t \sim \mathbf {c} \int_ {\beta} ^ {x} g (t) d t),
\]

for any \(\alpha\) and \(\beta\) in \([a, +\infty[\) .

It is enough to prove the proposition for the relation \(\mathbf{f} \ll g\) since, if \(\mathbf{c} \neq 0\) , the relation \(\mathbf{f} \sim \mathbf{c}g\) is equivalent to \(\mathbf{f} - \mathbf{c}g \ll g\) .

The first part is an immediate consequence of the mean value theorem, since if \(\|\mathbf{f}(x)\| \leqslant \varepsilon g(x)\) for \(x \geqslant x_{0}\) one deduces that

\[
\left\| \int_ {x} ^ {+ \infty} \mathbf {f} (t) d t \right\| \leqslant \int_ {x} ^ {+ \infty} \| \mathbf {f} (t) \| d t \leqslant \varepsilon \int_ {x} ^ {+ \infty} g (t) d t \quad \text { for } x \geqslant x _ {0}.
\]

In the second place, suppose that \(\int_{a}^{+\infty} g(t) dt = +\infty\) . If \(\| \mathbf{f}(x) \| \leqslant \varepsilon g(x)\) for \(x \geqslant x_0 \geqslant \max(\alpha, \beta)\) , one has

\[
\begin{array}{l} \int_ {\alpha} ^ {x} \| \mathbf {f} (t) \| d t = \int_ {\alpha} ^ {x _ {0}} \| \mathbf {f} (t) \| d t + \int_ {x _ {0}} ^ {x} \| \mathbf {f} (t) \| d t \\ \leqslant \int_ {\alpha} ^ {x _ {0}} \| \mathbf {f} (t) \| d t + \varepsilon \int_ {x _ {0}} ^ {x} g (t) d t \\ = \varepsilon \int_ {\beta} ^ {x} g (t) d t + \left(\int_ {\alpha} ^ {x _ {0}} \| \mathbf {f} (t) \| d t - \varepsilon \int_ {\beta} ^ {x _ {0}} g (t) d t\right). \end{array}
\]

Now there exists an \(x_{1} \geqslant x_{0}\) such that for every \(x \geqslant x_{1}\)

\[
\left| \int_ {\alpha} ^ {x _ {0}} \| \mathbf {f} (t) \| d t - \varepsilon \int_ {\beta} ^ {x _ {0}} g (t) d t \right| \leqslant \varepsilon \int_ {\beta} ^ {x} g (t) d t
\]

whence, for \(x \geqslant x_1\)

\[
\left\| \int_ {\alpha} ^ {x} \mathbf {f} (t) d t \right\| \leqslant \int_ {\alpha} ^ {x} \| \mathbf {f} (t) \| d t \leqslant 2 \varepsilon \int_ {\beta} ^ {x} g (t) d t
\]

which completes the proof, since \(\varepsilon > 0\) is arbitrary.

In other words, one can integrate the two terms of a strong relation \(\mathbf{f} \ll g\) , \(\mathbf{f} \sim \mathbf{a}g\) , when \(g\) is positive on an interval \([a, +\infty[\) , and the relation persists between the primitives of the two terms, provided one takes care to integrate from \(x\) to \(+\infty\) if \(\int_{a}^{+\infty} g(t) dt\) converges, and from \(\alpha\) to \(x\) (for any \(\alpha\) in \([a, +\infty[\) ) in the opposite case.

Note that props. 1 (V, p.~228) and 6 (V, p.~230) remain valid when \(\mathfrak{F}\) is a base for the trace filter of the intervals \([t, +\infty[\) (where \(t > a\) ) on the complement of a countable set (cf.~I, p.~15, th. 2).

Examples. 1) On applying prop. 6 of V, p.~230, to the relation \(1 / x \ll x^{\alpha - 1}\) where \(\alpha > 0\) , one again obtains the relation \(\log x \ll x^{\alpha}\) for every \(\alpha > 0\) , which is equivalent to the relation \(y^{1 / \alpha} \ll e^{y}\) proved in III, p.~105.

\begin{enumerate}
\def\labelenumi{\arabic{enumi})}
\setcounter{enumi}{1}
\tightlist
\item
  We have \(\left(\frac{e^x}{x}\right)' = \frac{e^x}{x}\left(1 - \frac{1}{x}\right) \sim e^x / x\) ; since \(e^x / x\) tends to \(+\infty\) with \(x\) , one deduces from prop. 6 of V, p.~230, that \(\int_{1}^{x} \frac{e^t}{t} dt \sim e^x / x\) .
\end{enumerate}

Remark. When \(g\) is not assumed to remain \(\geqslant 0\) on an interval \([a, +\infty[\) (or to remain \(\leqslant 0\) on such an interval), and \(\int_{a}^{+\infty} g(t) dt\) is not convergent, the relation \(f \sim g\) does not necessarily imply that \(\int_{a}^{x} f(t) dt \sim \int_{a}^{x} g(t) dt\) , as is shown by the example where \(g(x) = \sin x\) and \(f(x) = \left(1 + \frac{\sin x}{x}\right) \sin x\) ; here

\pandocbounded{\includegraphics[keepaspectratio]{images/part002_ecaa53ed.jpg}}

\[
\int_ {n \pi} ^ {(n + 1) \pi} \frac {\sin^ {2} t}{t} d t \geqslant \frac {1}{(n + 1) \pi} \int_ {0} ^ {\pi} \sin^ {2} t d t \geqslant \frac {1}{2} \int_ {n + 1} ^ {n + 2} \frac {d t}{t},
\]

whence

\[
\int_ {\pi} ^ {n \pi} \frac {\sin^ {2} t}{t} d t \geqslant \frac {1}{2} \int_ {2} ^ {n + 1} \frac {d t}{t}
\]

and the integral \(\int_{1}^{+\infty} dt / t\) is infinite, though \(\int_{\frac{\pi}{2}}^{x} g(t) dt = -\cos x\) remains bounded (see V, p.~260, exerc. 4).

